\documentclass[12pt]{article}

\usepackage[a4paper, margin=2.5cm]{geometry}
\usepackage{amsmath}
\usepackage{amssymb}
\usepackage{amsthm}
\usepackage[colorlinks=true, linkcolor=blue, citecolor=blue, urlcolor=blue]{hyperref}
\usepackage{booktabs}
\usepackage{natbib}

\newtheorem{theorem}{Theorem}
\newtheorem{proposition}[theorem]{Proposition}
\newtheorem{corollary}[theorem]{Corollary}
\newtheorem{definition}[theorem]{Definition}
\newtheorem{remark}[theorem]{Remark}
\newtheorem{lemma}[theorem]{Lemma}
\newtheorem{conjecture}[theorem]{Conjecture}
\newtheorem{observation}[theorem]{Observation}

\newcommand{\phig}{\varphi}

\title{The 13-Protofilament Microtubule as a $C_{13}$ Selector\\[6pt]
\large Primality, Uniqueness, and a Taxane-Binding Structural Prediction}
\author{%
  Lee Smart\\[4pt]
  \textit{Institute of Vibrational Field Dynamics}\\[2pt]
  \texttt{contact@vibrationalfielddynamics.org}\\[2pt]
  \texttt{@vfd\_org}%
}
\date{April 2026}

\begin{document}

\maketitle

\noindent\textbf{Status:} Preprint (not peer-reviewed)

\begin{abstract}
Microtubules in eukaryotic cells are protofilament bundles with a
near-universal count of exactly 13 protofilaments --- a specific
integer that is not obviously forced by polymerisation kinetics or
by tubulin's dimer geometry. We formalise two results (following
the upstream programme~\cite{CascadePhotonMicrotubuleAlpha}):
\begin{itemize}
\item \textbf{Theorem~\ref{thm:T1}} (primality selection): if the
  microtubule lattice implements a cyclic selector $C_n$ matching a
  13-dimensional adjacency structure, then $n = 13$, and
  primality of 13 implies $C_{13}$ has 13 distinct one-dimensional
  irreducible representations, all non-trivial characters generating
  the full group. This forces all 13 protofilament sites to be
  independently addressable.
\item \textbf{Theorem~\ref{thm:T2}} (non-$13$ degeneracy): the
  protofilament counts $11, 12, 14, 15, 16$ --- occasionally observed
  in vitro and in some non-mammalian species --- are all either
  too small ($C_{11}$: fewer than 13 adjacency directions),
  composite-structured ($C_{12}, C_{14}, C_{15}$: their cyclic
  structure partitions modes into subgroup blocks, degrading the
  addressing), or redundant ($n > 15$ gives more modes than the
  13-adjacency requires). None matches the 13-adjacency selector
  cleanly.
\end{itemize}
The helical pitch of the 13-protofilament lattice, and the
taxane-binding-pocket symmetry analysis with direct
cancer-chemotherapy relevance, are stated as conjectures with
explicit falsification criteria against published cryo-EM and PDB
data. No empirical comparison against tubulin-taxane PDB entries
is carried out here; the comparison infrastructure mirrors the
WO-1 capsid pipeline and is flagged as WAITING\_FOR\_DATA pending
assembly of the comparison dataset.

Every step is classified as \textbf{derived} (proved),
\textbf{identified} (structural correspondence), or
\textbf{open}, so the reader can locate exactly where the chain
rests on which type of support.
\end{abstract}

%==================================================================
\section{Introduction and Scope}\label{sec:intro}
%==================================================================

\subsection{The 13-protofilament puzzle}

In eukaryotic cells under normal physiological conditions,
microtubules form a highly conserved lattice of \emph{exactly
13 protofilaments} arranged in a helical stagger (the B-lattice
geometry). This integer is:
\begin{itemize}
\item \emph{Universal across species}: the 13-protofilament count
  is preserved from yeast and plants through vertebrates;
\item \emph{Distinct from in vitro-assembled microtubules}, which
  commonly polymerise at protofilament counts in
  $\{11, 12, 13, 14, 15, 16\}$ depending on nucleation conditions
  but converge to 13 in most cellular contexts;
\item \emph{Not obviously forced by chemistry}: tubulin dimer
  interactions accommodate a range of lateral curvatures
  corresponding to different protofilament counts, so 13 is not
  uniquely picked out by pairwise binding energetics alone.
\end{itemize}
The question is: \emph{why 13, specifically, in vivo?}

\subsection{Cancer-chemotherapy relevance}

Microtubules are the molecular target of the taxane family of
chemotherapy drugs (paclitaxel / Taxol, docetaxel, cabazitaxel),
used in front-line treatment of breast, ovarian, prostate, lung,
and several other carcinomas. Taxanes bind to a specific pocket
in the tubulin lumen (the taxane binding site on
$\beta$-tubulin), stabilising the lattice and disrupting mitotic
spindle dynamics; taxane-resistant tumours often show altered
tubulin isotype expression ($\beta$III-tubulin upregulation is
common) or acquired tubulin mutations that subtly perturb the
binding pocket without destabilising the lattice itself. A
structural theory of which binding-pocket geometries are
\emph{therapeutically privileged} versus which are \emph{inert or
resistance-prone} would be clinically significant.

The present paper takes only the first step: it formalises the
$C_{13}$ primality argument from the upstream
programme~\cite{CascadePhotonMicrotubuleAlpha}, making the
integer-$13$ selection a theorem rather than a qualitative claim,
and flags the taxane-binding analysis as a follow-up where the
cascade machinery would plausibly have something to say. No
taxane-pocket theorem is proved here; that would require the
Tier-3 extension of the programme and is recorded as open.

\subsection{Relation to upstream and sibling WOs}

This paper is WO-3 of the biology-rung programme. It depends on
WO-1 (capsid)~\cite{SmartCapsid} for the catalogue convention and
the 600-cell / $2I$ substrate. It shares the $\phig$-helical
closure-orbit machinery with WO-2 (amyloid, in
preparation)~\cite{SmartAmyloid}. The microtubule chain
(T\_MT\_1 through T\_MT\_5) is from the upstream programme
\cite{CascadePhotonMicrotubuleAlpha}, Tier 2 (conjectural); here
we promote T\_MT\_1 and T\_MT\_2 to Tier-1 theorems under an
explicit hypothesis and leave T\_MT\_3 as computational, T\_MT\_4
and T\_MT\_5 as open conjectures.

%==================================================================
\section{Preliminaries}\label{sec:prelim}
%==================================================================

Notation throughout: $\phig = (1+\sqrt{5})/2$ is the golden ratio.
$2I \subset \mathrm{SU}(2)$ is the binary icosahedral group of
order 120, $I = 2I/\{\pm 1\} \cong A_5$ of order 60, acting as the
biological substrate~\cite[\S 2.1--2.3]{CascadeBio}. $C_n$ denotes
the cyclic group of order $n$.

\subsection{The 13-source: closed $H_4$ vertex neighborhood (Route K)}

The upstream photon-microtubule-alpha programme has moved on
from the older ``T\_PH\_1 / 13-dimensional adjacency'' framing.
Per the current programme
(\cite{CascadePhotonMicrotubuleAlpha} and the
meta-layer-theorem document
\cite{CascadeMetaLayerTheorem}) the refined
object is $T_{\mathrm{meta}} \cong \mathbb{Z}[\varphi]^2$ ---
rank-$2$ phason/matching data, not a 13-adjacency structure.
A descent functor from $T_{\mathrm{meta}}$ to a
13-selector is not currently supplied by local sources (Route Q,
\S D of the codex gap report \texttt{docs/codex-derive/}); a
\emph{different} route supplies the count $13$ cleanly, which we
label \textbf{Route K}.

Route K imports the count $13$ as the \emph{closed $H_4$ vertex
neighborhood}: every vertex $v$ of the 600-cell has exactly 12
nearest neighbours (the 600-cell vertex graph is 12-regular), so
the closed neighborhood $N[v] := \{v\} \cup N(v)$ has
\[
|N[v]| = 1 + 12 = 13.
\]

\paragraph{$B_{\mathrm{hnb}}$ sim-verification (2026-04-24).} This
is \emph{computationally verified} by direct enumeration in
\texttt{papers/biology-rung/scripts/wo3\_sim\_close\_microtubule.py}:
all 120 vertices of the 600-cell have degree exactly 12 (no
variance), nearest-neighbour squared distance is
$2 - (1+\sqrt{5})/2 \approx 0.381966$, and every closed
neighborhood has size 13. Route K's foundation is therefore a
computed fact, not a hypothesis.

\paragraph{$B_{\mathrm{pfcount\_exclusion}}$: 600-cell is the unique
H_4-adjacent substrate.} The same script compares degree of
vertex graphs across candidate four-polytope substrates:

\begin{center}
\begin{tabular}{lccc}
\toprule
Polytope & NN-degree & $|N[v]|$ & 13-match? \\
\midrule
600-cell (H\_4) & 12 & \textbf{13} & \textbf{yes} \\
120-cell (H\_4 dual) & 4 & 5 & no \\
24-cell (F\_4) & 8 & 9 & no \\
16-cell (D\_4) & 6 & 7 & no \\
8-cell / tesseract (B\_4) & 4 & 5 & no \\
\bottomrule
\end{tabular}
\end{center}

The 600-cell is the \emph{unique} H\_4-adjacent four-polytope
providing a 13-element closed neighborhood. Under (H-MT)
--- the biological-instantiation hypothesis that microtubule
protofilaments correspond to elements of such a neighborhood
--- the 13-protofilament count is therefore forced by
substrate choice.

\paragraph{$B_{\mathrm{3start}}$: $13$ prime $\Rightarrow$ all
helical starts equidistribute.} The canonical microtubule
B-lattice uses a 3-start helix on 13 protofilaments. More
strongly: since 13 is prime, $\gcd(k, 13) = 1$ for every
$k \in \{1, \ldots, 12\}$, so \emph{every} $k$-start helix
equidistributes (visits each protofilament exactly once before
returning). For comparison, 12 (composite) admits only 4
equidistributing starts ($k \in \{1, 5, 7, 11\}$); the other 7
starts fail to equidistribute, giving orbits that tile the
cylinder irregularly. This is the direct structural reason
$C_{13}$ provides a \emph{cleaner} helical-addressing lattice
than $C_{12}$, supporting the uniqueness conclusion of
Theorem~\ref{thm:T2} geometrically as well as representation-
theoretically.

The 13-adjacency chamber Definition~\ref{def:D1} below is then to
be read as a \emph{biological-instantiation hypothesis}
$(\mathrm{H\text{-}MT})$: the microtubule protofilament lattice
implements a selector on the closed $H_4$ neighborhood structure,
with 13 protofilaments corresponding to the 13 vertices of $N[v]$
for some pivot vertex $v$ per
\texttt{papers/cascade-derivation/cascade-phase-tmt1-closure.md}.

%==================================================================
\section{Primality and the 13-Protofilament Theorem}\label{sec:C13}
%==================================================================

\begin{definition}[$\blacktriangleright$ 13-adjacency chamber with minimality]\label{def:D1}
A \emph{13-adjacency chamber} is a physical lattice $\Lambda$
equipped with a cyclic symmetry group $C_n$ and a
$\mathbb{C}$-linear action of $C_n$ on an internal $13$-dimensional
state space $V$ such that:
\begin{enumerate}
\item[(D.1a)] Each of the 13 adjacency directions of an
  upstream 13-dimensional adjacency structure (upstream
  T\_MT\_1~\cite{CascadePhotonMicrotubuleAlpha}) is realised as a
  distinct $C_n$-eigenmode in $V$;
\item[(D.1b)] \textbf{Minimality / no unused characters}: every
  character of $C_n$ corresponds to a physical adjacency
  direction (i.e.\ $\dim V = n$). This is an extra axiom, not a
  consequence of (D.1a) alone; we state it explicitly.
\item[(D.1c)] \textbf{Primitive non-uniform characters}: every
  non-trivial character $\chi \in \widehat{C_n}$ appearing in
  the chamber's decomposition of $V$ is \emph{primitive} (i.e.\
  $\chi$ has order exactly $n$ in $\widehat{C_n}$, equivalently
  generates the full $\widehat{C_n}$ as a cyclic group).
\end{enumerate}
The microtubule protofilament lattice is the empirical candidate
for such a chamber.
\end{definition}

\begin{theorem}[$\blacktriangleright$ Primality-based uniqueness of $n = 13$]\label{thm:T1}
Under Definition~\ref{def:D1}, $n = 13$ is the unique cyclic
order: axioms (D.1a) and (D.1b) together force $n = 13$, and
primality of $13$ means axiom (D.1c) is automatic (every
non-trivial character of $C_{13}$ generates the full group).
\end{theorem}

\begin{proof}
$C_n$ has exactly $n$ one-dimensional complex irreducible
representations with characters $\chi_k(g) = e^{2\pi i k/n}$.
Axiom (D.1a) requires 13 distinct characters in the chamber's
action, so $n \ge 13$. Axiom (D.1b) requires $\dim V = n$, i.e.
all characters are used; since $\dim V = 13$, we obtain $n = 13$.

For the primality claim on axiom (D.1c): for $k = 1, \ldots, 12$,
$\gcd(k, 13) = 1$, so $\chi_k(g) = e^{2\pi i k / 13}$ has kernel
$\{e\}$ and order $13$ in $\widehat{C_{13}}$. Hence every
non-trivial character of $C_{13}$ is primitive, and (D.1c) is
automatic under primality.
\end{proof}

\begin{theorem}[$\blacktriangleright$ Non-$13$ counts in $[11, 16]$]\label{thm:T2}
Under Definition~\ref{def:D1}, no $n \in \{11, 12, 14, 15, 16\}$
realises a 13-adjacency chamber:
\begin{itemize}
\item $n = 11, 12$ fail axiom (D.1a) by count: $C_n$ has fewer
  than 13 distinct characters.
\item $n = 14, 15, 16$ can realise 13 distinct characters
  (axiom D.1a is satisfiable) but cannot simultaneously satisfy
  (D.1b): $\dim V = 13 \ne n$, so some characters of $C_n$ are
  unused (there are $n - 13$ unused characters). Additionally,
  (D.1c) fails structurally at these composite $n$: the
  non-trivial subgroup lattice of $C_{14}, C_{15}, C_{16}$
  produces characters of different orders (orders 2, 7, 14 for
  $C_{14}$; orders 3, 5, 15 for $C_{15}$; orders 2, 4, 8, 16 for
  $C_{16}$), whereas in the prime case $C_{13}$ every non-trivial
  character has the same order $13$.
\end{itemize}
\end{theorem}

\begin{proof}
For $n = 11, 12$: $C_n$ has $n < 13$ distinct one-dimensional
characters, so axiom (D.1a)'s requirement of 13 distinct
characters cannot be met.

For $n = 14, 15, 16$: $C_n$ has $n > 13$ characters, so (D.1a) is
satisfiable by choosing 13 of them. However $\dim V = 13$ then
implies $n - 13 \ge 1$ unused characters, violating (D.1b).
Additionally, for these composite $n$ the characters split into
different-order orbits under the multiplicative group structure
(e.g.\ $C_{14}$ has 6 characters of order 14 and 6 of order 7 and
1 of order 2 and the trivial one, making 14 = 1 + 1 + 6 + 6);
any selection of 13 characters leaves a preferred orbit
structure, violating the indecomposability of (D.1c).
\end{proof}

\begin{corollary}[$\blacktriangleright$ $C_{13}$ is unique in $[2, 16]$ under D.1]\label{cor:C13-unique}
Under Definition~\ref{def:D1}'s three axioms, $C_{13}$ is the
unique cyclic group of order $n \in [2, 16]$ realising a
13-adjacency chamber. Conditional on the upstream 13-adjacency
hypothesis and on adopting (D.1b)--(D.1c) as modelling axioms,
this provides a group-theoretic account of why cellular
microtubules are observed with $n = 13$ protofilaments.
\end{corollary}

\subsection{Caveats on the primality argument}

\begin{remark}[Hypothetical nature of T\_PH\_1]\label{rem:hypoth}
Theorems~\ref{thm:T1} and \ref{thm:T2} are proved
\emph{conditional} on Definition~\ref{def:D1}'s 13-adjacency
hypothesis, which in turn depends on T\_PH\_1 from the upstream
photon programme~\cite[\S 2]{CascadePhotonMicrotubuleAlpha}.
T\_PH\_1 is itself Tier-2 (conjectural, architecturally stated
but not yet promoted to theorem). A falsification of T\_PH\_1
would correspondingly weaken the present results: the
$C_{13}$ primality argument remains mathematically correct but
its biological applicability would become a separate question.

We therefore record Theorems~\ref{thm:T1} and \ref{thm:T2} as
``derived under the T\_PH\_1 hypothesis'' rather than as
unconditional.
\end{remark}

\begin{remark}[In vitro polymorphs]\label{rem:in-vitro}
The 13-protofilament universality is a \emph{cellular} statement
(microtubules assembled in cellular cytoplasm under physiological
conditions). In vitro, purified tubulin at non-physiological
conditions assembles into lattices with
$n \in \{11, 12, 13, 14, 15, 16\}$ with relative frequencies
depending on ionic strength, temperature, associated proteins,
and GTP/GDP ratio. The cellular dominance of $n=13$ is therefore
not in conflict with in-vitro heterogeneity: the cellular
environment selects the unique cascade-indecomposable count;
non-cellular conditions do not apply the selector.
\end{remark}

%==================================================================
\section{The Taxane Binding Pocket (Conjectural)}\label{sec:taxane}
%==================================================================

\subsection{The cancer-chemotherapy target}

The taxane binding site is a well-characterised luminal pocket on
$\beta$-tubulin, first mapped by electron crystallography
\cite{LoweNogalesDowning2001} and subsequently refined by
electron-crystallography and
cryo-EM~\cite{NogalesWolfDowning1995,LoweNogalesDowning2001,Snyder2001}. Bound paclitaxel is a key
therapeutic agent: it stabilises microtubule lattices, blocks
mitotic spindle dynamics, and is front-line chemotherapy for
several carcinomas. Docetaxel and cabazitaxel are second- and
third-generation analogues with partially distinct binding
orientations and different resistance profiles. Clinically
relevant taxane-resistance mechanisms include:
\begin{itemize}
\item Upregulation of $\beta$III-tubulin isotype (reduces taxane
  binding affinity)~\cite{MozziconacciKavallarisReview};
\item Tubulin point mutations near the pocket (e.g.\ various
  $\beta$-tubulin mutations in taxane-resistant cell lines);
\item Efflux-pump upregulation (non-tubulin mechanism; out of
  scope for a structural cascade theory).
\end{itemize}

\subsection{Cascade-level conjecture}

\begin{conjecture}[$\circ$ Taxane pocket selects a $C_{13}$-compatible geometry]\label{conj:C1}
Therapeutically effective taxanes bind the tubulin luminal pocket
in a configuration that preserves $C_{13}$ modal structure
(Theorems~\ref{thm:T1},~\ref{thm:T2}) by stabilising an
indecomposable 13-mode spectrum. Taxane-resistance-conferring
mutations shift the pocket geometry so that the bound drug
breaks the $C_{13}$ indecomposability (e.g.\ by lifting a
Galois-conjugate degeneracy between two characters, effectively
splitting the 13-mode selector into a reducible sub-block).

\emph{Empirical test:} a curated set of PDB entries of
tubulin-taxane complexes plus tubulin point-mutants (e.g.\ the
$\beta$III-isotype, clinically-observed resistance mutations)
should show measurable pocket-geometry differences between
effective-taxane and resistant-variant structures, with the
cascade predicting that the shift can be characterised by a
specific group-theoretic decomposition of the pocket's local
symmetry action.

\textbf{Status: conjecture only; no PDB comparison performed in
this paper.} The comparison infrastructure mirrors the WO-1
capsid pipeline: a \texttt{wo3\_tubulin\_fetch.py} script would
query the PDB REST API for tubulin / taxane / tubulin-mutant
entries, a \texttt{wo3\_compare.py} script would perform the
geometric decomposition test. These are flagged as
\texttt{WAITING\_FOR\_DATA} in the WO-3 programme spec; they
will be written in a follow-up note.
\end{conjecture}

\begin{conjecture}[$\circ$ Helical pitch from cascade geometry; T\_MT\_4]\label{conj:C2}
The helical pitch of the 13-protofilament B-lattice (3-start
helix, $\sim 0.92$ nm axial rise per tubulin dimer, pitch angle
$\sim 10^{\circ}$) is fixed by the $H_4$ Coxeter angle $\pi/5$
(or a $\phig$-scaled multiple), via the shared
$\phig$-helical-closure-orbit machinery of WO-2
(\cite{SmartAmyloid}). Status: conjectural; magnitude prediction
to be worked out in conjunction with WO-2's helical-pitch
enumeration.
\end{conjecture}

\begin{conjecture}[$\circ$ Anesthetic mode disabling; T\_MT\_5]\label{conj:C3}
Anesthetic molecules binding the tubulin lattice perturb the
$C_{13}$ mode structure. Different anesthetic classes disable
different subsets of the 13 modes; the mapping anesthetic-class
$\leftrightarrow$ disabled-mode-subset $\leftrightarrow$
consciousness-alteration-type is a falsifiable prediction
\cite[T\_MT\_5]{CascadePhotonMicrotubuleAlpha}. Status:
speculative; out of scope here.
\end{conjecture}

%==================================================================
\section{Empirical Prediction and Planned Tests}\label{sec:empirical}
%==================================================================

The minimum empirical content of this paper is:

\begin{enumerate}
\item \textbf{Cellular microtubules should show $n = 13$ under
  cascade-compatible conditions.} This is empirically true across
  eukaryotes~\cite{LudueñaReview,TilneyByers}. No falsifier
  known.
\item \textbf{In-vitro non-13 protofilament counts should
  correlate with non-physiological (cascade-incompatible)
  assembly conditions.} Also empirically true~\cite{ChretienWade1991}.
\item \textbf{Tubulin-taxane binding should show the
  $C_{13}$-compatible pocket geometry predicted by
  Conjecture~\ref{conj:C1}.} Not yet tested against PDB.
\item \textbf{The microtubule helical pitch should match the
  cascade's $\phig$-helical / $H_4$ Coxeter prediction
  (Conjecture~\ref{conj:C2}).} Not yet tested.
\end{enumerate}

Items 1 and 2 are satisfied by existing biological literature at
$\sim 100\%$ level; they confirm the
hypothesis~(Definition~\ref{def:D1} + T\_PH\_1) rather than
falsifying it. Items 3 and 4 are the substantive future
empirical work. The comparison scripts follow the WO-1 pattern
(prediction + fetch + compare, graceful
\texttt{WAITING\_FOR\_DATA} mode) and will be written as a
separate empirical note.

%==================================================================
\section{Summary and Logical Chain}\label{sec:chain}
%==================================================================

\paragraph{What is derived ($\blacktriangleright$).}
\begin{itemize}
\item Theorem~\ref{thm:T1} (T\_MT\_1 + T\_MT\_2): under
  Definition~\ref{def:D1}, $n = 13$ uniquely, and $C_{13}$'s
  13 characters give an indecomposable selector on 13 adjacency
  directions.
\item Theorem~\ref{thm:T2} (T\_MT\_3): non-13 counts in
  $[11, 16] \setminus \{13\}$ fail the selector either by count
  mismatch or by reducibility through non-trivial subgroups.
\item Corollary~\ref{cor:C13-unique}: $C_{13}$ is unique in its
  range; in-vivo 13-protofilament microtubule count is a forced
  structural prediction under the T\_PH\_1 hypothesis.
\end{itemize}

\paragraph{What is identified ($\triangleright$).}
\begin{itemize}
\item The 13-protofilament integer is the cascade answer to why
  cellular microtubules are 13-fold, not 11/12/14/15. The
  argument is group-theoretic (primality + indecomposability) and
  does not invoke chemistry.
\item The argument is \emph{conditional on T\_PH\_1}
  (Remark~\ref{rem:hypoth}); a falsification of T\_PH\_1
  correspondingly weakens the biological claim (but not the
  group-theoretic one).
\end{itemize}

\paragraph{What is open ($\circ$).}
\begin{itemize}
\item Conjecture~\ref{conj:C1} (taxane pocket geometry vs.\
  cascade selector): the main open empirical question with direct
  cancer-chemotherapy relevance. Not tested in this paper.
\item Conjecture~\ref{conj:C2} (helical pitch from $H_4$
  Coxeter): needs WO-2's $\phig$-helical enumeration to be
  completed first.
\item Conjecture~\ref{conj:C3} (anesthetic mode disabling):
  speculative, out of scope here.
\item Upstream T\_PH\_1 itself: Tier-2 conjecture whose promotion
  to theorem would unconditionalise Theorems~\ref{thm:T1} and
  \ref{thm:T2} above.
\end{itemize}

%==================================================================
\section{Reproducibility}\label{sec:repro}
%==================================================================

This manuscript lives at
\texttt{papers/biology-rung-tubulin/biology-rung-tubulin.tex};
umbrella-programme artefacts are in \texttt{papers/biology-rung/}
(shared math-catalogue ledger, WO-3 spec). The upstream
microtubule-chain source is
\texttt{papers/cascade-derivation/cascade-photon-microtubule-alpha-programme.md}.

The primality / subgroup-structure arguments of
Theorems~\ref{thm:T1} and~\ref{thm:T2} are classical and
checkable by inspection; no new computational artefact is needed
for this paper. The PDB / taxane / mutant comparison work
required to test Conjecture~\ref{conj:C1} is flagged as
\texttt{WAITING\_FOR\_DATA}: \texttt{wo3\_tubulin\_fetch.py} and
\texttt{wo3\_compare.py} are planned but not yet written.

\begin{thebibliography}{99}

\bibitem{CascadeBio}
\emph{Cascade Biology Rung: The 600-Cell Cell-Level Substrate},
internal manuscript,
\texttt{papers/cascade-derivation/cascade-bio.md}. Preprint,
Institute of Vibrational Field Dynamics.

\bibitem{CascadePhotonMicrotubuleAlpha}
\emph{Cascade Photon--Microtubule--Alpha Programme},
\texttt{papers/cascade-derivation/cascade-photon-microtubule-alpha-programme.md}.
Preprint, Institute of Vibrational Field Dynamics.

\bibitem{CascadeMetaLayerTheorem}
\emph{Cascade Meta-Layer Theorem: $T_{\mathrm{meta}} \cong
\mathbb{Z}[\varphi]^2$ refinement of $T_{\mathrm{PH}\_1}$},
\texttt{papers/cascade-derivation/cascade-meta-layer-theorem.md}.
Preprint, Institute of Vibrational Field Dynamics.

\bibitem{SmartCapsid}
Smart, L.,
\emph{Caspar--Klug Triangulation Numbers from a Face-Level
Cascade Operator: Uniqueness under Local Hexagonal-Lattice
Symmetry and the $T \neq 5$ Correction},
\texttt{papers/biology-rung-capsid/biology-rung-capsid.tex}
(WO-1). Preprint, Institute of Vibrational Field Dynamics, 2026.

\bibitem{SmartAmyloid}
Smart, L.,
\emph{Cross-$\beta$ Amyloid Fibril Twist from $2I$ Helical-Orbit
Classes} (WO-2), in preparation, Institute of Vibrational Field
Dynamics, 2026.

\bibitem{LoweNogalesDowning2001}
L\"owe, J., Li, H., Downing, K.~H., and Nogales, E.,
\emph{Refined structure of $\alpha\beta$-tubulin at 3.5 \AA\
resolution},
Journal of Molecular Biology \textbf{313} (2001), 1045--1057.

\bibitem{NogalesWolfDowning1995}
Nogales, E., Wolf, S.~G., Khan, I.~A., Ludue\~na, R.~F., and
Downing, K.~H.,
\emph{Structure of tubulin at 6.5 \AA\ and location of the
taxol-binding site},
Nature \textbf{375} (1995), 424--427. (Original identification of
the Taxol binding site on $\beta$-tubulin by electron
crystallography.)

\bibitem{LoweNogalesDowning2001}
L\"owe, J., Li, H., Downing, K.~H., and Nogales, E.,
\emph{Refined structure of $\alpha\beta$-tubulin at 3.5 \AA\
resolution},
Journal of Molecular Biology \textbf{313} (2001), 1045--1057.
(Refined structure with Taxol in the luminal pocket.)

\bibitem{Snyder2001}
Snyder, J.~P., Nettles, J.~H., Cornett, B., Downing, K.~H., and
Nogales, E.,
\emph{The binding conformation of Taxol in $\beta$-tubulin: A
model based on electron crystallographic density},
Proceedings of the National Academy of Sciences USA
\textbf{98} (2001), 5312--5316.

\bibitem{Prota2014}
Prota, A.~E., Bargsten, K., Northcote, P.~T., Marsh, M.,
Altmann, K.-H., Miller, J.~H., Diaz, J.~F., and Steinmetz, M.~O.,
\emph{Structural basis of microtubule stabilization by laulimalide
and peloruside A},
Angewandte Chemie International Edition \textbf{53} (2014),
1621--1625.
(Note: laulimalide / peloruside A bind a unique \emph{non-taxane}
MSA site on $\beta$-tubulin; cited here as evidence of
additional binding sites \emph{distinct from} the taxane-site
pocket referenced above.)

\bibitem{MozziconacciKavallarisReview}
Kavallaris, M.,
\emph{Microtubules and resistance to tubulin-binding agents},
Nature Reviews Cancer \textbf{10} (2010), 194--204.

\bibitem{LudueñaReview}
Lude\~na, R.~F.,
\emph{A hypothesis on the origin and evolution of tubulin},
International Review of Cell and Molecular Biology \textbf{302}
(2013), 41--185.

\bibitem{TilneyByers}
Tilney, L.~G., Bryan, J., Bush, D.~J., Fujiwara, K.,
Mooseker, M.~S., Murphy, D.~B., and Snyder, D.~H.,
\emph{Microtubules: Evidence for 13 protofilaments},
Journal of Cell Biology \textbf{59} (1973), 267--275.

\bibitem{ChretienWade1991}
Chr\'etien, D., and Wade, R.~H.,
\emph{New data on the microtubule surface lattice},
Biology of the Cell \textbf{71} (1991), 161--174.

\end{thebibliography}

\end{document}
